\section*{Capítulo \ref{ch:b3:orogenic-wedges} --- Cinturones de cabalgamientos y cuñas orogénicas}

\begin{solution}{exo:b3:orogenic-wedges:1}
$\sin\ang{30} = 0.5$: $K = 2\times0.5/0.5 = 2$ y $\sigma'_1/\sigma'_3 = 1 +
K = 3$.
\end{solution}

\begin{solution}{exo:b3:orogenic-wedges:2}
$\phi = \arctan0.6 = \ang{31.0}$, $\sin\phi = 0.514$, $K = 2.12$. En una
base horizontal, $\alpha = 0.3/3.12 = \qty{0.0962}{rad} = \ang{5.5}$. En una
base que buza $\ang{3} = \qty{0.0524}{rad}$, $\alpha + \beta = 0.352/3.12 =
\qty{0.113}{rad} = \ang{6.5}$, así que $\alpha = \ang{3.5}$.
\end{solution}

\begin{solution}{exo:b3:orogenic-wedges:3}
Profundidad $= (T - 10)/30$: apatito (U--Th)/He, de 1.3 a \qty{2.3}{km}; traza de fisión
en apatito, de 2.7 a \qty{3.7}{km}; traza de fisión en circón $270/30 = \qty{9.0}{km}$. A
\qty{5}{km}, la roca estaba a $10 + 150 = \qty{160}{\celsius}$: circón
(U--Th)/He, que cierra de 160 a \qty{200}{\celsius}, a profundidades de 5.0 a
\qty{6.3}{km}.
\end{solution}

\begin{solution}{exo:b3:orogenic-wedges:4}
Longitud horizontal $2/\tan\ang{30} = \qty{3.46}{km}$; a lo largo del buzamiento
$2/\sin\ang{30} = \qty{4.0}{km}$. El estrato se eleva en la altura de la rampa,
\qty{2}{km}. La línea de corte debe recorrer toda la rampa: un desplazamiento de
\qty{4.0}{km}.
\end{solution}

\begin{solution}{exo:b3:orogenic-wedges:5}
$\phi = \arctan1.03 = \ang{45.8}$, $K = 5.08$; $\beta = \qty{0.105}{rad}$.
Para $\lambda = 0.4$: $\alpha + \beta = (0.105 + 0.85\times0.6)/(1 +
5.08\times0.6) = 0.615/4.05 = \qty{0.152}{rad} = \ang{8.7}$, $\alpha =
\ang{2.7}$. Para $\lambda = 0.9$: $0.190/1.51 = \qty{0.126}{rad} = \ang{7.2}$,
$\alpha = \ang{1.2}$. El aumento de la presión de fluidos reduce el \omterm{def:b3:orogenic-wedges:critical-taper}{ahusamiento crítico}:
la cuña se vuelve supercrítica, se desliza hacia delante sobre su base sin
engrosarse, y su frente avanza hasta que la acreción y la erosión la devuelven
al nuevo ahusamiento.
\end{solution}

\begin{solution}{exo:b3:orogenic-wedges:6}
$S = A/r = 6.0/1.5 = \qty{4.0}{km}$. Si se interpreta como \omterm{def:b3:orogenic-wedges:thrust-folds}{pliegue de despegue}, $S = A/h =
6.0/3.0 = \qty{2.0}{km}$: la mitad del valor real. Las dos interpretaciones deben distinguirse
mediante la forma del pliegue y la presencia de una rampa.
\end{solution}

\begin{solution}{exo:b3:orogenic-wedges:7}
$K = 2.12$ y $1 + K(1 - 0.4) = 2.27$. La ecuación da $\mu_b(1 -
\lambda_b) = (\alpha + \beta)\times2.27 - \beta$, con $\beta =
\qty{0.0349}{rad}$. Jura: $0.0611\times2.27 - 0.0349 = 0.10$. Sur:
$0.218\times2.27 - 0.0349 = 0.46$. Byerlee a presión hidrostática:
$0.85\times0.6 = 0.51$. La base del sur está cerca de la de Byerlee; la
base del Jura es casi cinco veces más débil, como cabe esperar de un despegue en
halita, cuya fricción es baja, bajo fluidos que no están sobrepresurizados.
\end{solution}

\begin{solution}{exo:b3:orogenic-wedges:8}
$\kappa = \qty{e-6}{m^2/s} = \qty{31.6}{km^2/Myr}$, y $v =
\qty{1}{km/Myr}$: $\mathrm{Pe} = 30/31.6 = 0.951$. El factor $0.951/(1 -
e^{-0.951}) = 1.55$ eleva el gradiente de 30 a \qty{46.5}{\celsius/km}.
Para \qty{65}{\celsius}: $1 - e^{-\mathrm{Pe}\,z/L} = (55/900)(1 - e^{-0.951})
= 0.0375$, así que $\mathrm{Pe}\,z/L = 0.0382$ y $z = \qty{1.21}{km}$, frente a
$55/30 = \qty{1.83}{km}$ sin erosión.
\end{solution}

\begin{solution}{exo:b3:orogenic-wedges:9}
$\mathrm dh/\mathrm dt = \qty{0.8}{km}/\qty{1.2}{Myr} = \qty{0.67}{mm/yr}$.
Con $f = 0.5$, $v = 0.67\times(1 - 0.5) = \qty{0.33}{mm/yr}$.
\end{solution}

\begin{solution}{exo:b3:orogenic-wedges:10}
Con $\alpha = 0$ y $u = 1 - \lambda$: $\beta(1 + Ku) = \beta + \mu_b u$,
así que $\beta Ku = \mu_b u$ y $\beta^* = \mu_b/K$, sea cual sea $u$. Para Taiwán,
$0.85/5.08 = \qty{0.167}{rad} = \ang{9.6}$. Sobre una base más empinada, la pendiente crítica
de la superficie es negativa: cualquier cuña con una superficie horizontal o inclinada hacia el antepaís
ya es supercrítica, y se desliza sobre su base sin engrosarse. Un
cinturón de cabalgamientos no puede adquirir allí su forma normal.
\end{solution}

\begin{solution}{exo:b3:orogenic-wedges:11}
Cuando $\Lambda \to \infty$ y $v \to 0$, $m \to 0$, así que $e^{-mz_c} \to 1$, y
la geoterma es recta, $G_0 = G(z_c) = G$: $f \to (G - \Gamma)/G = 1 -
\Gamma/G = 1 - 6.5/30 = 0.78$. La pendiente es entonces $v/(1 - f) = vG/\Gamma =
4.6v$. Para $v = \qty{1}{mm/yr}$, $v/2\kappa = \qty{0.0158}{km^{-1}}$. A
\qty{5}{km}: $m = 0.0158 + \sqrt{0.0158^2 + 1.257^2} = \qty{1.273}{km^{-1}}$,
$e^{-1.273\times1.21} = 0.215$, $f = 40.0\times0.215/44.7 = 0.19$. A
\qty{20}{km}: $m = 0.0158 + 0.3146 = \qty{0.330}{km^{-1}}$, $e^{-0.399} =
0.671$, $f = 40.0\times0.671/44.7 = 0.60$.
\end{solution}

\begin{solution}{exo:b3:orogenic-wedges:12}
La sección es un triángulo de base $W$ y altura $W\tan(\alpha + \beta)$:
$A = W^2\tan(\alpha + \beta)/2$. Con $\mathrm dA/\mathrm dt = F_a$ y $A =
0$ en $t = 0$, $A = F_at$, lo que da $W$. Con $\tan\ang{6} = 0.105$:
$W = \sqrt{2\times40\times2/0.105} = \qty{39}{km}$ tras \qty{2}{Myr}; $t =
60^2\times0.105/(2\times40) = \qty{4.7}{Myr}$ para alcanzar \qty{60}{km}.
\end{solution}

\begin{solution}{pb:b3:orogenic-wedges:1}
\textbf{1.} $4.5 - 1.5 = \qty{3.0}{km}$ de altura; $3.0/\tan\ang{25} =
\qty{6.43}{km}$ de longitud horizontal; $3.0/\sin\ang{25} = \qty{7.10}{km}$ a lo largo
del buzamiento.

\textbf{2.} Por la altura de la rampa, \qty{3.0}{km}.

\textbf{3.} $S = A/r = 16.5/3.0 = \qty{5.5}{km}$.

\textbf{4.} Lo mismo, \qty{16.5}{km^2}: en un \omterm{def:b3:orogenic-wedges:thrust-folds}{pliegue de flexión de falla} el exceso de área
de cada estrato sobre el rellano superior es $Sr$. Como \omterm{def:b3:orogenic-wedges:thrust-folds}{pliegue de despegue}, $16.5/4.0 =
\qty{4.1}{km}$: un cuarto menos de lo que corresponde.

\textbf{5.} $4\times(6.0 - 2.5) = \qty{14.0}{km}$, dos veces y media la
\qty{5.5}{km} de desplazamiento en la rampa frontal.

\textbf{6.} El caballo frontal, el más cercano al antepaís: cada rampa nueva
se forma en el bloque inferior, por delante de la anterior. Cada caballo, elevado por su rampa,
eleva el cabalgamiento de techo y la cobertera situada encima, como un \omterm{def:b3:orogenic-wedges:thrust-folds}{pliegue de flexión de falla} eleva su
bloque superior, de modo que el techo se pliega sobre el apilamiento.

\textbf{7.} $\alpha = \arctan(2.1/60) = \ang{2.00}$; $\beta =
\arctan(4.2/60) = \ang{4.00}$; ahusamiento \ang{6.0}.

\textbf{8.} $2.4 + 8.2 = \qty{10.6}{km}$, menos de unos \qty{15}{km}: la
base sigue siendo friccional, y se aplica la teoría de Coulomb.

\textbf{9.} $\phi = \ang{45.8}$, $\sin\phi = 0.717$, $K = 5.08$. En seco:
$\alpha + \beta = (0.0699 + 0.85)/6.08 = \qty{0.151}{rad} = \ang{8.7}$,
$\alpha = \ang{4.7}$.

\textbf{10.} $\lambda = 1000/2500 = 0.40$: $(0.0699 + 0.51)/4.05 =
\qty{0.143}{rad} = \ang{8.2}$, $\alpha = \ang{4.2}$. El valor observado \ang{2.0}
haría que la cuña fuera fuertemente subcrítica, engrosándose en lugar de avanzar:
los fluidos deben estar sobrepresurizados.

\textbf{11.} Con $T = \qty{0.1049}{rad}$ y $u = 1 - \lambda$, $T(1 + Ku) =
\beta + \mu_bu$ da $u = (\beta - T)/(TK - \mu_b) = (0.0699 -
0.1049)/(0.5327 - 0.85) = 0.110$: $\lambda = 0.89$.

\textbf{12.} $\sigma_v = 2500\times9.81\times10\,600 = \qty{260}{MPa}$; $P =
0.89\times260 = \qty{231}{MPa}$; hidrostática $1000\times9.81\times10\,600 =
\qty{104}{MPa}$. Los fluidos soportan \qty{127}{MPa} más de lo que soportaría
una columna de agua.

\textbf{13.} $(0.0699 + 0.85\times0.3)/(1 + 5.08\times0.3) = 0.325/2.52 =
\qty{0.129}{rad} = \ang{7.4}$, $\alpha = \ang{3.4}$. La cuña, con \ang{2.0},
es ahora subcrítica: deja de avanzar y se engrosa mediante \omterm{def:b3:orogenic-wedges:thrust-geometry}{cabalgamientos fuera de
secuencia} detrás de su frente hasta que la pendiente de su superficie alcance \ang{3.4}.

\textbf{14.} $F_a = \qty{10}{km/Myr}\times\qty{4.0}{km} = \qty{40}{km^2/Myr}$,
es decir, \qty{4.0e4}{m^3/yr} por kilómetro de cinturón.

\textbf{15.} $\bar e = 40/60 = \qty{0.67}{km/Myr} = \qty{0.67}{mm/yr}$.

\textbf{16.} $40\times100 = \qty{4000}{km^3/Myr} = \qty{4.0e6}{m^3/yr}$, o
$\qty{1.0e10}{kg/yr}$ (10 millones de toneladas al año).

\textbf{17.} $\mathrm dW/\mathrm dt = F_a/(W\tan(\alpha + \beta)) =
40/(60\times0.105) = \qty{6.3}{km/Myr}$.

\textbf{18.} No: en estado estacionario la erosión debe seguir eliminando todo el
flujo entrante, de modo que la tasa media se mantiene en \qty{0.67}{mm/yr}. Pero la roca añadida
bajo la retaguardia se eleva allí: la erosión y la exhumación serían más rápidas
sobre la retaguardia, donde se encontrarían las edades de enfriamiento más jóvenes.

\textbf{19.} $v = \qty{0.667}{km/Myr}$: $\mathrm{Pe} = 0.667\times30/31.6 =
0.634$; $G_0 = 30\times0.634/(1 - e^{-0.634}) = 30\times1.35 =
\qty{40.5}{\celsius/km}$.

\textbf{20.} $1 - e^{-\mathrm{Pe}\,z/L} = (55/900)(1 - e^{-0.634}) = 0.0287$,
luego $\mathrm{Pe}\,z/L = 0.0291$ y $z_c = \qty{1.38}{km}$; $G(z_c) =
40.5\,e^{-0.0291} = \qty{39.3}{\celsius/km}$.

\textbf{21.} $1.38/0.667 = \qty{2.07}{Myr}$, la edad de la muestra a
\qty{2.00}{km}: a la altitud media la deflexión de la isoterma es nula.

\textbf{22.} $v/2\kappa = 0.667/63.1 = \qty{0.0106}{km^{-1}}$, $2\pi/\Lambda =
\qty{0.314}{km^{-1}}$: $m = 0.0106 + 0.3143 = \qty{0.325}{km^{-1}}$;
$e^{-0.325\times1.38} = 0.639$; $f = (40.5 - 6.5)\times0.639/39.3 = 0.552$.

\textbf{23.} La altitud media es \qty{2.00}{km} y la edad media
\qty{2.068}{Myr}; $\sum\Delta t\,\Delta h = 0.823$ y $\sum\Delta t^2 =
0.552$, por tanto la pendiente es \qty{1.49}{km/Myr}: una tasa aparente de
\qty{1.49}{mm/yr}.

\textbf{24.} $v = 1.49\times(1 - 0.552) = \qty{0.667}{mm/yr}$.

\textbf{25.} El \omterm{def:b3:orogenic-wedges:thermochronology}{termocronómetro} y el balance de flujo concuerdan: la cuña está
en estado estacionario, y tomar la pendiente aparente tal cual habría
sobrestimado la tasa más del doble. La tasa de exhumación de la cuña Kestrel
es \textbf{\qty[text-series-to-math=true]{0.67}{mm/yr}}.
\end{solution}
